\chapter{Phân tích Yêu cầu}
\label{chap:requirements}

\section{Yêu cầu Người dùng (User Requirements)}
\label{sec:req-user}

Đối tượng người dùng chính là sinh viên ngành Cơ Kỹ Thuật đang thực hiện
bài tập lớn thiết kế hệ dẫn động. Qua trao đổi với nhóm CKT, các yêu cầu
tổng quan được xác định như sau:

\begin{itemize}
    \item Người dùng muốn nhập các thông số đầu bài (công suất, số vòng quay,
            thời gian phục vụ, phương án) và nhận kết quả tính toán ngay lập tức.
    \item Người dùng muốn xem gợi ý động cơ phù hợp thay vì phải tra bảng thủ công.
    %   \item Người dùng muốn \textbf{thay đổi một thông số} và xem toàn bộ kết quả
            cập nhật tức thì mà không cần tính lại từ đầu.
    \item Người dùng muốn lưu lại kết quả để tham khảo khi viết báo cáo.
    \item Ứng dụng phải dùng được không cần kết nối Internet
            (phòng thi, thư viện không có WiFi).
\end{itemize}

\section{Yêu cầu Chức năng (Functional Requirements)}
\label{sec:req-functional}

\subsection{Module 1 -- Tính Chọn Động cơ}

\begin{itemize}
  \item \textbf{FR1.1}: Người dùng nhập các thông số đầu vào:
        công suất trục công tác $P_t$~(kW), số vòng quay $n$~(v/p),
        thời gian phục vụ $L$~(năm), chế độ tải.
  \item \textbf{FR1.2}: Hệ thống tự động tính hiệu suất $\eta$ từ bảng cố định
        (bộ truyền côn--trụ--xích), công suất cần thiết $P_{ct}$
        và số vòng quay sơ bộ $n_{sb}$.
  \item \textbf{FR1.3}: Hệ thống tra cứu CSDL và hiển thị danh sách động cơ
        thỏa điều kiện $P_{dc} \geq P_{ct}$ và $n_{db} \approx n_{sb}$;
        người dùng chọn động cơ phù hợp.
  \item \textbf{FR1.4}: Hệ thống tính lại tỉ số truyền thực $u_t$, phân phối
        $u_h$ và $u_x$, xuất bảng công suất -- mô-men -- số vòng quay
        trên từng trục (I, II, III, công tác).
\end{itemize}

\subsection{Module 2 -- Tính Bộ truyền Xích}

\begin{itemize}
  \item \textbf{FR2.1}: Module nhận tự động tham số trục xích
        ($P$, $n$, $T$) từ kết quả Module~1.
  \item \textbf{FR2.2}: Hệ thống tính chọn số răng đĩa xích $z_1, z_2$;
        bước xích $p_c$; khoảng cách trục $a$; số mắc xích $X$.
  \item \textbf{FR2.3}: Hệ thống kiểm nghiệm: số lần va đập $i \leq [i]$;
        hệ số an toàn $s \geq [s]$; độ bền tiếp xúc $\sigma_H \leq [\sigma_H]$;
        áp suất $p \leq [p]$.
  \item \textbf{FR2.4}: Hiển thị bảng tổng kết thông số bộ truyền xích,
        kết quả kiểm nghiệm (đạt / không đạt).
\end{itemize}

\subsection{Module 3 -- Tính Bộ truyền Bánh răng}

\begin{itemize}
  \item \textbf{FR3.1}: Module nhận tự động tham số từng trục từ kết quả Module~1.
  \item \textbf{FR3.2}: \textbf{Cấp nhanh -- bánh răng côn}: chọn vật liệu và ứng suất
        cho phép; xác định thông số hình học (chiều dài côn $R_e$, số răng $z_1, z_2$,
        mô-đun $m_{te}$, góc côn $\delta_1, \delta_2$); kiểm nghiệm tiếp xúc và uốn.
  \item \textbf{FR3.4}: Hiển thị bảng tổng kết đầy đủ thông số bánh răng.
\end{itemize}

\subsection{Quản lý Phiên tính toán}

\begin{itemize}
  \item \textbf{FR4.1}: Người dùng có thể lưu phiên tính toán hiện tại với tên tự đặt.
  \item \textbf{FR4.2}: Người dùng có thể tải lại phiên đã lưu và tiếp tục chỉnh sửa.
  \item \textbf{FR4.3}: Người dùng có thể xóa phiên đã lưu.
\end{itemize}

\section{Yêu cầu Phi chức năng (Non-Functional Requirements)}
\label{sec:req-nonfunctional}

\begin{itemize}
  \item \textbf{NFR1 -- Hiệu năng}: Toàn bộ kết quả tính toán hiển thị trong vòng
        1~giây sau khi người dùng nhấn \textit{Tính toán}.
  \item \textbf{NFR2 -- Độ chính xác}: Kết quả các giá trị số (công suất, mô-men,
        thông số hình học) sai số không quá $\pm 0{,}5\%$ so với tính tay
        theo giáo trình~\cite{trinh2006tinh}.
  \item \textbf{NFR3 -- Khả năng sử dụng}: Sinh viên chưa quen lập trình có thể
        hoàn thành tính toán Module~1 mà không cần đọc hướng dẫn.
  \item \textbf{NFR4 -- Offline}: Toàn bộ tính năng hoạt động không cần kết nối mạng;
        bảng tra cứu được đóng gói sẵn trong ứng dụng.
  \item \textbf{NFR5 -- Nền tảng}: Hỗ trợ Android~10+ và iOS~14+.
\end{itemize}
\section{Use Case Diagram}
\label{sec:req-usecase}

Ứng dụng có một tác nhân chính là \textbf{Sinh viên CKT} (người dùng cuối).
Các use case chính được xác định như sau:

\begin{itemize}
  \item \textbf{UC01} -- Tính chọn động cơ điện và phân phối tỉ số truyền.
  \item \textbf{UC02} -- Tính thiết kế và kiểm nghiệm bộ truyền xích (\texttt{<<include>>} UC01).
  \item \textbf{UC03} -- Tính thiết kế và kiểm nghiệm bộ truyền bánh răng (\texttt{<<include>>} UC01).
  \item \textbf{UC04} -- Lưu và tải lại phiên tính toán.
  \item \textbf{UC05} -- Đăng nhập / Đăng ký tài khoản.
\end{itemize}

UC02 và UC03 có quan hệ \texttt{<<include>>} với UC01:
kết quả tính chọn động cơ từ UC01 (công suất, số vòng quay, mô-men trên các trục)
là đầu vào bắt buộc cho cả hai module tiếp theo.

\begin{figure}[H]
\centering
\begin{tikzpicture}[scale=0.88,
  uc/.style={ellipse, draw=black, minimum width=4.8cm, minimum height=1.05cm,
             align=center, font=\small},
  inc/.style={dashed, ->, >=stealth}
]
  %--- System boundary ---
  \draw[thick, rounded corners=5pt] (0.5,-0.5) rectangle (12.0,14.0);
  \node[font=\bfseries\small] at (6.25,13.55) {MechCalc Mobile};

  %--- Use Cases ---
  \node[uc] (uc05) at (6.25,12.2) {UC05 \\ Đăng nhập / Đăng ký};
  \node[uc] (uc01) at (6.25, 9.4) {UC01 \\ Tính chọn động cơ};
  \node[uc] (uc02) at (3.3,  6.0) {UC02 \\ Tính bộ truyền xích};
  \node[uc] (uc03) at (9.2,  6.0) {UC03 \\ Tính bánh răng};
  \node[uc] (uc04) at (6.25, 2.6) {UC04 \\ Lưu / Tải phiên};

  %--- Stick-figure actor at (-2.1, 7.2) ---
  \draw[fill=white] (-2.1,9.0) circle (0.38);   % head
  \draw (-2.1,8.62) -- (-2.1,7.7);              % body
  \draw (-2.6,8.15) -- (-1.6,8.15);             % arms
  \draw (-2.1,7.7)  -- (-2.5,6.85);             % left leg
  \draw (-2.1,7.7)  -- (-1.7,6.85);             % right leg
  \node[font=\small, align=center] at (-2.1,6.35) {Sinh viên\\CKT};

  %--- Associations (actor right-arm tip → each UC) ---
  \draw (-1.6,8.15) -- (uc05.west);
  \draw (-1.6,8.15) -- (uc01.west);
  \draw (-1.6,8.15) -- (uc02.west);
  \draw (-1.6,8.15) -- (uc03.west);
  \draw (-1.6,8.15) -- (uc04.west);

  %--- <<include>> stereotypes ---
  \draw[inc] (uc02.north east) -- (uc01.south west)
        node[midway, above left, font=\scriptsize]
        {$\langle\!\langle$include$\rangle\!\rangle$};
  \draw[inc] (uc03.north west) -- (uc01.south east)
        node[midway, above right, font=\scriptsize]
        {$\langle\!\langle$include$\rangle\!\rangle$};
\end{tikzpicture}
\caption{Use Case Diagram -- Ứng dụng MechCalc Mobile}
\label{fig:usecase-diagram}
\end{figure}

\section{Đặc tả Use Case Chi tiết}
\label{sec:req-usecase-detail}

\subsection{UC01 -- Tính chọn động cơ}

\begin{table}[H]
\centering
\caption{Đặc tả UC01 -- Tính chọn động cơ}
\label{tab:uc01}
\begin{tabularx}{\textwidth}{|l|X|}
\hline
\textbf{Tác nhân} & Sinh viên CKT \\
\hline
\textbf{Điều kiện tiên quyết} & Người dùng đã đăng nhập thành công. \\
\hline
\textbf{Luồng chính} &
  1. Người dùng nhập $P_{lv}$~(kW), $n_{lv}$~(v/p), $L$~(năm). \newline
  2. Hệ thống tính hiệu suất $\eta$, công suất cần thiết $P_{ct}$ và số vòng quay sơ bộ $n_{sb}$. \newline
  3. Hệ thống truy vấn CSDL, hiển thị danh sách tối đa 5 động cơ gợi ý (sắp xếp theo độ gần $n_{đb} \approx n_{sb}$). \newline
  4. Người dùng chọn một động cơ phù hợp. \newline
  5. Hệ thống tính tỉ số truyền thực tế $u_t$, $u_h$, $u_1$, $u_2$ và lập bảng $P$--$n$--$T$ trên 4 trục. \\
\hline
\textbf{Luồng thay thế} &
  3a. Không có động cơ nào thỏa $P_{đc} \geq P_{ct}$ → hiển thị thông báo ``Không tìm thấy động cơ phù hợp''. \\
\hline
\textbf{Hậu điều kiện} &
  Kết quả Module~1 được lưu vào \texttt{globalNavigationState} để Module~2 và 3 sử dụng. \\
\hline
\end{tabularx}
\end{table}

\subsection{UC02 -- Tính bộ truyền xích}

\begin{table}[H]
\centering
\caption{Đặc tả UC02 -- Tính bộ truyền xích}
\label{tab:uc02}
\begin{tabularx}{\textwidth}{|l|X|}
\hline
\textbf{Tác nhân} & Sinh viên CKT \\
\hline
\textbf{Điều kiện tiên quyết} & UC01 đã hoàn thành (tồn tại $P_{III}$, $n_{III}$, $u_x$ trong \texttt{globalNavigationState}). \\
\hline
\textbf{Luồng chính} &
  1. Hệ thống tự động nạp $P$, $n$, $u$ từ kết quả Module~1. \newline
  2. Hệ thống tính số răng $z_1$, $z_2$; bước xích $p_c$; khoảng cách trục $a$; số mắt xích $X$. \newline
  3. Hệ thống thực hiện 4 kiểm nghiệm: số lần va đập $i$; hệ số an toàn $s$; ứng suất tiếp xúc $\sigma_H$; áp suất $p$. \newline
  4. Hệ thống hiển thị bảng tổng kết kết quả với trạng thái đạt/không đạt từng điều kiện. \\
\hline
\textbf{Hậu điều kiện} &
  Kết quả bộ truyền xích được lưu vào \texttt{globalNavigationState.chainResult}. \\
\hline
\end{tabularx}
\end{table}

\subsection{UC03 -- Tính bộ truyền bánh răng}

\begin{table}[H]
\centering
\caption{Đặc tả UC03 -- Tính bộ truyền bánh răng}
\label{tab:uc03}
\begin{tabularx}{\textwidth}{|l|X|}
\hline
\textbf{Tác nhân} & Sinh viên CKT \\
\hline
\textbf{Điều kiện tiên quyết} & UC01 đã hoàn thành (tồn tại $T_I$, $n_I$, $u_1$, $T_{II}$, $n_{II}$, $u_2$). \\
\hline
\textbf{Luồng chính} &
  1. Hệ thống tự động nạp tham số từng trục từ kết quả Module~1. \newline
  2. \textbf{Cấp nhanh (bánh răng côn)}: chọn vật liệu, tính kích thước sơ bộ $R_e$, xác định $z_1$, $z_2$, $m_{te}$, $\delta_1$, $\delta_2$; kiểm nghiệm $\sigma_H$ và $\sigma_F$. \newline
  3. \textbf{Cấp chậm (bánh răng trụ)}: tính khoảng cách trục sơ bộ $a_w$, chọn mô-đun $m$, xác định $z_1$, $z_2$, $b_w$; kiểm nghiệm $\sigma_H$ và $\sigma_F$. \newline
  4. Hệ thống hiển thị bảng tổng kết đầy đủ thông số hình học và kết quả kiểm nghiệm. \\
\hline
\textbf{Hậu điều kiện} &
  Kết quả bánh răng được lưu vào \texttt{globalNavigationState.gearResult}. \\
\hline
\end{tabularx}
\end{table}

\subsection{UC04 -- Lưu và tải phiên tính toán}

\begin{table}[H]
\centering
\caption{Đặc tả UC04 -- Lưu và tải phiên tính toán}
\label{tab:uc04}
\begin{tabularx}{\textwidth}{|l|X|}
\hline
\textbf{Tác nhân} & Sinh viên CKT \\
\hline
\textbf{Điều kiện tiên quyết} & Người dùng đã đăng nhập. \\
\hline
\textbf{Luồng chính -- Lưu} &
  1. Người dùng nhấn nút \textit{Lưu dự án} và nhập tên phiên. \newline
  2. Hệ thống gọi API backend lưu toàn bộ thông số đầu vào và kết quả tính toán. \newline
  3. Phiên mới xuất hiện trong danh sách tab \textit{Dự án}. \\
\hline
\textbf{Luồng chính -- Tải} &
  1. Người dùng chọn một phiên từ danh sách Dự án. \newline
  2. Hệ thống nạp kết quả vào \texttt{globalNavigationState}. \newline
  3. Người dùng có thể xem lại kết quả của từng module. \\
\hline
\textbf{Luồng chính -- Xóa} &
  1. Người dùng vuốt hoặc nhấn nút Xóa trên một phiên. \newline
  2. Hệ thống gọi API xóa phiên và cập nhật danh sách. \\
\hline
\end{tabularx}
\end{table}

% ============================================================
\section{Sơ đồ Tuần tự (Sequence Diagram)}
\label{sec:req-sequence}

\subsection{SD01 -- UC01: Tính chọn động cơ}

\begin{figure}[H]
\centering
\resizebox{\textwidth}{!}{%
\begin{tikzpicture}[
  part/.style={rectangle, draw=black, minimum width=3.0cm, minimum height=0.8cm,
               align=center, font=\small},
  msg/.style={ ->, >=stealth, font=\scriptsize},
  ret/.style={dashed, ->, >=stealth, font=\scriptsize},
  life/.style={dashed, gray, thin}
]
  %--- Participants ---
  \node[part] (U)  at ( 0,  0) {Sinh viên};
  \node[part] (UI) at ( 5,  0) {\texttt{calc.tsx}};
  \node[part] (SV) at (11,  0) {\small\texttt{motor-dk-selection.ts}};
  \node[part] (DB) at (16,  0) {SQLite DB};

  %--- Lifelines ---
  \draw[life] ( 0,-0.4) -- ( 0,-13.5);
  \draw[life] ( 5,-0.4) -- ( 5,-13.5);
  \draw[life] (11,-0.4) -- (11,-13.5);
  \draw[life] (16,-0.4) -- (16,-13.5);

  %--- Messages ---
  \draw[msg] ( 0,-1.3) -- ( 5,-1.3)
        node[midway,above] {Nhập $P_{lv}$, $n_{lv}$, $L$};

  \draw[msg] ( 5,-2.6) -- (11,-2.6)
        node[midway,above] {tính $\eta$, $P_{ct}$, $n_{sb}$};
  \draw[ret] (11,-3.5) -- ( 5,-3.5)
        node[midway,above] {$P_{ct} = 9{,}06$ kW,\quad $n_{sb} = 2100$ v/p};

  \draw[msg] ( 5,-4.8) -- (16,-4.8)
        node[midway,above] {\texttt{SELECT WHERE power\_kw} $\geq P_{ct}$};
  \draw[ret] (16,-5.8) -- ( 5,-5.8)
        node[midway,above] {danh sách 5 ứng viên động cơ};

  \draw[msg] ( 5,-7.0) -- ( 0,-7.0)
        node[midway,above] {Hiển thị danh sách gợi ý};

  \draw[msg] ( 0,-8.3) -- ( 5,-8.3)
        node[midway,above] {Chọn IE3-W41R 160 M2 (11 kW, 2950 v/p)};

  \draw[msg] ( 5,-9.5) -- (11,-9.5)
        node[midway,above] {phân phối $u$, tính bảng $P$-$n$-$T$};
  \draw[ret] (11,-10.4) -- ( 5,-10.4)
        node[midway,above] {$u_1=4{,}7$;\quad $u_2=4{,}18$;\quad bảng $P$-$n$-$T$ (4 trục)};

  \draw[msg] ( 5,-11.7) -- ( 0,-11.7)
        node[midway,above] {Hiển thị bảng kết quả};

  % Self-note: write to globalNavigationState
  \draw[ret] ( 5,-12.8) to[out=0,in=0] ( 5,-13.0)
        node[right=4pt, font=\scriptsize] {ghi kết quả vào \texttt{globalNavigationState}};
\end{tikzpicture}
}
\caption{Sequence Diagram -- UC01: Tính chọn động cơ}
\label{fig:sd-uc01}
\end{figure}

\subsection{SD02 -- UC04: Lưu phiên tính toán}

\begin{figure}[H]
\centering
\resizebox{0.92\textwidth}{!}{%
\begin{tikzpicture}[
  part/.style={rectangle, draw=black, minimum width=3.0cm, minimum height=0.8cm,
               align=center, font=\small},
  msg/.style={ ->, >=stealth, font=\scriptsize},
  ret/.style={dashed, ->, >=stealth, font=\scriptsize},
  life/.style={dashed, gray, thin}
]
  %--- Participants ---
  \node[part] (U)  at ( 0, 0) {Sinh viên};
  \node[part] (UI) at ( 5, 0) {\texttt{project.tsx}};
  \node[part] (AX) at (10, 0) {Axios Service};
  \node[part] (BE) at (15, 0) {Spring Backend};

  %--- Lifelines ---
  \draw[life] ( 0,-0.4) -- ( 0,-9.5);
  \draw[life] ( 5,-0.4) -- ( 5,-9.5);
  \draw[life] (10,-0.4) -- (10,-9.5);
  \draw[life] (15,-0.4) -- (15,-9.5);

  %--- Messages ---
  \draw[msg] ( 0,-1.3) -- ( 5,-1.3)
        node[midway,above] {Nhấn Lưu, nhập tên phiên};

  \draw[msg] ( 5,-2.6) -- (10,-2.6)
        node[midway,above] {\texttt{saveProject(name, data)}};

  \draw[msg] (10,-3.8) -- (15,-3.8)
        node[midway,above] {POST /api/projects};

  \draw[ret] (15,-4.9) -- (10,-4.9)
        node[midway,above] {201 Created \{id, name\}};

  \draw[ret] (10,-6.1) -- ( 5,-6.1)
        node[midway,above] {project saved};

  \draw[msg] ( 5,-7.3) -- ( 0,-7.3)
        node[midway,above] {Cập nhật danh sách Dự án};

  % Alt: network error
  \draw[thick] (0,-8.3) -- (15,-8.3);
  \node[font=\scriptsize, anchor=west] at (0,-8.05) {[alt] Lỗi mạng};
  \draw[ret] (10,-8.9) -- ( 5,-8.9)
        node[midway,above] {Network Error → hiển thị thông báo lỗi};
\end{tikzpicture}
}
\caption{Sequence Diagram -- UC04: Lưu phiên tính toán}
\label{fig:sd-uc04}
\end{figure}
